% !TeX spellcheck = en_US
% !TeX encoding = UTF-8
\chapter*{Abstract}

This thesis presents the design and implementation of a vulnerability scanner that detects security flaws in Python applications using two complementary approaches: traditional static code analysis and modern machine learning techniques. The primary objective is to identify common web application vulnerabilities, including SQL injection, Cross-Site Scripting (XSS), command injection, and Cross-Site Request Forgery (CSRF).

The proposed vulnerability scanner is integrated into a Python Flask-based backend with a React frontend, providing a user-friendly interface for analyzing and comparing results from both detection methods. The system allows users to upload source code files, paste code directly into the interface, or scan public GitHub repositories. By leveraging both rule-based static analysis and data-driven machine learning models, the tool offers a comprehensive assessment of potential security risks.

This solution is particularly beneficial for programmers and developers, enabling early detection of vulnerabilities during the development lifecycle and contributing to improved software security practices.

This vulnerability scanner is opensource and is available at: https://github.com/mohammadalmasi/06.27

\textbf{Index Terms—}Vulnerability Scanner, Static Analysis, Machine Learning techniques.
